% GENERATED FILE. Edit canonical lessons, then run npm run generate:books.
% canonical-source-hash: fnv1a64:ccf39ab01d86be84
% canonical-lessons: KA-C158-beyisu, KA-C158-aduge-madu, KA-C158-kharidisu, KA-C158-pavatisu, KA-C158-snana-madu

\chapter{To boil, To cook, To buy, To pay, To bathe}
\label{ch:ka-to-boil-to-cook-to-buy-to-pay-to-bathe}
\hlchaptermodality{158}

\begin{chapteropening}
\textbf{By the end of this chapter:} \emph{I can say to boil, to cook, to buy, to pay and to bathe.}

Complete the last lesson of chapter 158: I can say to boil, to cook, to buy, to pay and to bathe.
\end{chapteropening}

\section[\texorpdfstring{\kn{ಬೇಯಿಸು} (bēyisu)}{ಬೇಯಿಸು (bēyisu)}]{\kn{ಬೇಯಿಸು} (bēyisu) — to boil, to cook (in water or steam)}

\label{lesson:KA-C158-beyisu}



\begin{quote}
Before the new one: say the Kannada for an advertisement, then the Kannada for a receipt.
\end{quote}

\subsection*{You'll want to know: \kn{ಬೇಯಿಸು}}
\textbf{\kn{ಬೇಯಿಸು}} — \emph{bēyisu} — \textquotedblleft{}to boil, to cook (in water or steam)\textquotedblright{}.

\begin{grammarlens}[title={What you've built}]
One more word for this chapter.
\end{grammarlens}

\subsection*{Your turn}
\begin{itemize}
\raggedright
  \item \emph{Say it:} \emph{bēyisu}
  \item \emph{Say it:} \emph{bēyisu}, once more
  \item \emph{Say it:} say \emph{bēyisu}
  \item \emph{Recall:} say the Kannada for an advertisement, then the Kannada for a receipt, then say \emph{bēyisu} again
\end{itemize}

\begin{tcolorbox}[breakable,colback=teal!4,colframe=teal!35!black,title={Before you move on}]
What does \kn{ಬೇಯಿಸು} mean? (\textquotedblleft{}To boil, to cook (in water or steam)\textquotedblright{}.) Say it once more.
\end{tcolorbox}
\section[\texorpdfstring{\kn{ಅಡುಗೆ} \kn{ಮಾಡು} (aḍuge māḍu)}{ಅಡುಗೆ ಮಾಡು (aḍuge māḍu)}]{\kn{ಅಡುಗೆ} \kn{ಮಾಡು} (aḍuge māḍu) — to cook}

\label{lesson:KA-C158-aduge-madu}



\begin{quote}
Before the new one: say the Kannada for a receipt, then the Kannada for to boil.
\end{quote}

\subsection*{You'll want to know: \kn{ಅಡುಗೆ} \kn{ಮಾಡು}}
\textbf{\kn{ಅಡುಗೆ} \kn{ಮಾಡು}} — \emph{aḍuge māḍu} — \textquotedblleft{}to cook\textquotedblright{}.

\begin{grammarlens}[title={What you've built}]
One more word for this chapter.
\end{grammarlens}

\subsection*{Your turn}
\begin{itemize}
\raggedright
  \item \emph{Say it:} \emph{aḍuge māḍu}
  \item \emph{Say it:} \emph{aḍuge māḍu}, once more
  \item \emph{Say it:} say \emph{aḍuge māḍu}
  \item \emph{Recall:} say the Kannada for a receipt, then the Kannada for to boil, then say \emph{aḍuge māḍu} again
\end{itemize}

\begin{tcolorbox}[breakable,colback=teal!4,colframe=teal!35!black,title={Before you move on}]
What does \kn{ಅಡುಗೆ} \kn{ಮಾಡು} mean? (\textquotedblleft{}To cook\textquotedblright{}.) Say it once more.
\end{tcolorbox}
\section[\texorpdfstring{\kn{ಖರೀದಿಸು} (kharīdisu)}{ಖರೀದಿಸು (kharīdisu)}]{\kn{ಖರೀದಿಸು} (kharīdisu) — to buy}

\label{lesson:KA-C158-kharidisu}



\begin{quote}
Before the new one: say the Kannada for to boil, then the Kannada for to cook.
\end{quote}

\subsection*{You'll want to know: \kn{ಖರೀದಿಸು}}
\textbf{\kn{ಖರೀದಿಸು}} — \emph{kharīdisu} — \textquotedblleft{}to buy\textquotedblright{}.

\begin{grammarlens}[title={What you've built}]
One more word for this chapter.
\end{grammarlens}

\subsection*{Your turn}
\begin{itemize}
\raggedright
  \item \emph{Say it:} \emph{kharīdisu}
  \item \emph{Say it:} \emph{kharīdisu}, once more
  \item \emph{Say it:} say \emph{kharīdisu}
  \item \emph{Recall:} say the Kannada for to boil, then the Kannada for to cook, then say \emph{kharīdisu} again
\end{itemize}

\begin{tcolorbox}[breakable,colback=teal!4,colframe=teal!35!black,title={Before you move on}]
What does \kn{ಖರೀದಿಸು} mean? (\textquotedblleft{}To buy\textquotedblright{}.) Say it once more.
\end{tcolorbox}
\section[\texorpdfstring{\kn{ಪಾವತಿಸು} (pāvatisu)}{ಪಾವತಿಸು (pāvatisu)}]{\kn{ಪಾವತಿಸು} (pāvatisu) — to pay}

\label{lesson:KA-C158-pavatisu}



\begin{quote}
Before the new one: say the Kannada for to cook, then the Kannada for to buy.
\end{quote}

\subsection*{You'll want to know: \kn{ಪಾವತಿಸು}}
\textbf{\kn{ಪಾವತಿಸು}} — \emph{pāvatisu} — \textquotedblleft{}to pay\textquotedblright{}.

\begin{grammarlens}[title={What you've built}]
One more word for this chapter.
\end{grammarlens}

\subsection*{Your turn}
\begin{itemize}
\raggedright
  \item \emph{Say it:} \emph{pāvatisu}
  \item \emph{Say it:} \emph{pāvatisu}, once more
  \item \emph{Say it:} say \emph{pāvatisu}
  \item \emph{Recall:} say the Kannada for to cook, then the Kannada for to buy, then say \emph{pāvatisu} again
\end{itemize}

\begin{tcolorbox}[breakable,colback=teal!4,colframe=teal!35!black,title={Before you move on}]
What does \kn{ಪಾವತಿಸು} mean? (\textquotedblleft{}To pay\textquotedblright{}.) Say it once more.
\end{tcolorbox}
\section[\texorpdfstring{\kn{ಸ್ನಾನ} \kn{ಮಾಡು} (snāna māḍu)}{ಸ್ನಾನ ಮಾಡು (snāna māḍu)}]{\kn{ಸ್ನಾನ} \kn{ಮಾಡು} (snāna māḍu) — to bathe}

\label{lesson:KA-C158-snana-madu}



\begin{quote}
Before the new one: say the Kannada for to buy, then the Kannada for to pay.
\end{quote}

\subsection*{You'll want to know: \kn{ಸ್ನಾನ} \kn{ಮಾಡು}}
\textbf{\kn{ಸ್ನಾನ} \kn{ಮಾಡು}} — \emph{snāna māḍu} — \textquotedblleft{}to bathe\textquotedblright{}.

\begin{grammarlens}[title={What you've built}]
One more word for this chapter.
\end{grammarlens}

\subsection*{Your turn}
\begin{itemize}
\raggedright
  \item \emph{Say it:} \emph{snāna māḍu}
  \item \emph{Say it:} \emph{snāna māḍu}, once more
  \item \emph{Say it:} say \emph{snāna māḍu}
  \item \emph{Recall:} say the Kannada for to buy, then the Kannada for to pay, then say \emph{snāna māḍu} again
\end{itemize}

\begin{tcolorbox}[breakable,colback=teal!4,colframe=teal!35!black,title={Before you move on}]
What does \kn{ಸ್ನಾನ} \kn{ಮಾಡು} mean? (\textquotedblleft{}To bathe\textquotedblright{}.) Say it once more.
\end{tcolorbox}
